\section{Block-reset chains}\label{sec:blockreset}

The proof of Theorem \ref{thm:reduction} uses the reverse-erasure structure
only through two facts: a test function enters through its arithmetic and
geometric means, and a power sum is maximised at those means by a vector that
singles out a rarest atom. This section shows that both facts survive in a
much larger class of chains, with no restriction on the laziness. The
consequence is that every chain in the class has the same coefficient as a
small explicit chain, in which each block is replaced by a two-point law.

\subsection{The class and its collapse}

Let \(X\) be the disjoint union of finite nonempty blocks \(X_B\),
\(B\in\mathcal B\), and let \(\mu_B\) be a full-support law on \(X_B\). Fix
numbers \(a_B\in[0,1]\) and \(M_{BC}\ge0\) with
\(\sum_C M_{BC}=1-a_B\). The \emph{block-reset chain} is
\begin{equation}\label{br:eq:K}
K(x,\cdot)=a_B\,\delta_x+\sum_{C}M_{BC}\,\mu_C\qquad(x\in X_B).
\end{equation}
From \(x\) it stays put with probability \(a_B\); otherwise it chooses a block
\(C\) and resamples a fresh point from \(\mu_C\). The cross weights between a
point and a block are therefore rank one inside each pair of blocks.
Take a reference law \(\pi=\sum_B w_B\mu_B\) with \(w_B>0\),
\(\sum_B w_B=1\). Stationarity is not required. Then
\begin{equation}\label{br:eq:output}
\lambda:=\pi K=\sum_C w'_C\,\mu_C,\qquad
w'_C=a_Cw_C+\sum_B w_B M_{BC}.
\end{equation}
We assume \(w'_C>0\) for every \(C\), so that \(\lambda\) has full support;
this holds automatically when \(\pi K=\pi\).
Write \(q_B=\min\mu_B\). When \(|X_B|\ge2\), \(q_B\le1/2\).

The \emph{collapsed chain} \(\widehat K\) has the same block indices,
the same numbers \(a_B,M_{BC},w_B\), and block laws
\[
\widehat\mu_B=(q_B,\,1-q_B)\ \text{ on two points if }|X_B|\ge2,\qquad
\widehat\mu_B=(1)\ \text{ if }|X_B|=1.
\]
Its reference law is \(\widehat\pi=\sum_B w_B\widehat\mu_B\). It has at most
\(2|\mathcal B|\) states, whatever the sizes of the original blocks.

\begin{lemma}[Closure under composition]\label{br:lem:closure}
Let \(K\) and \(K'\) be block-reset chains on the same blocks and block laws,
with coefficients \((a,M)\) and \((a',M')\). Then \(KK'\) is a block-reset
chain on the same blocks and laws, with
\[
a''_B=a_Ba'_B,\qquad
M''_{BD}=a_BM'_{BD}+M_{BD}a'_D+\sum_C M_{BC}M'_{CD}.
\]
These coefficients do not involve the laws \(\mu_B\), so collapsing commutes
with composition: \(\widehat{KK'}=\widehat K\,\widehat{K'}\).
\end{lemma}
\begin{proof}
For \(x\in X_B\),
\(\mu_CK'=a'_C\mu_C+\sum_D M'_{CD}\mu_D\). Substituting this into
\(KK'(x,\cdot)=a_BK'(x,\cdot)+\sum_C M_{BC}\,\mu_CK'\) gives the stated
coefficients.
\end{proof}

\noindent\textit{In words.} Running two block-reset steps in a row is again a
block-reset step, so every power of a chain in the class is in the class.

\subsection{Power sums at fixed arithmetic and geometric means}

For \(0<w<1\), \(0<g<1\) and \(p>0\), let \(a<b\) be the two positive values
with \((1-w)a+wb=1\) and \(a^{1-w}b^w=g\), and put
\(\Psi_p(w)=(1-w)a^p+wb^p\), exactly as in Lemma \ref{lem:weight}.
(Existence and uniqueness of \(a<b\) is shown in the proof of that lemma
and does not use \(p\).) Here \(b\) is the larger value and \(w\) its
weight.

\begin{lemma}[Weight monotonicity for every exponent]\label{br:lem:weight}
Fix \(0<g<1\). The function \(w\mapsto \Psi_p(w)\) on \((0,1)\) is strictly
decreasing if \(p>1\), strictly increasing if \(0<p<1\), and constant equal
to \(1\) if \(p=1\).
\end{lemma}
\begin{proof}
The case \(p=1\) is the constraint itself. Let \(p\ne1\), \(r=b/a>1\), and
write primes for derivatives in \(w\) at fixed \(g\). Put
\(\xi_1=(1-w)a'\) and \(\xi_2=wb'\) (called \(A\) and \(B\) in Lemma
\ref{lem:weight}). Differentiating the two constraints gives
\[
\xi_1+\xi_2=a-b,\qquad \frac{\xi_1}a+\frac{\xi_2}b=-\log r .
\]
Solving this linear system,
\[
\xi_2=-\frac{b\,(r-1-\log r)}{r-1},\qquad
\xi_1=\frac{a\,(r-1-r\log r)}{r-1}.
\]
Hence
\[
\Psi_p'(w)=b^p-a^p+p\big(a^{p-1}\xi_1+b^{p-1}\xi_2\big)
=\frac{a^p}{r-1}\Big[p\log r\,(r^p-r)-(p-1)(r-1)(r^p-1)\Big],
\]
which is \eqref{eq:weightderivative}; the elimination holds for every real
\(p\). Put \(I_1=\int_1^r t^{p-1}dt\) and \(I_0=\int_1^r t^{p-2}dt\), so that
\(r^p-1=pI_1\) and \(r^p-r=(p-1)rI_0\). The bracket equals
\[
p(p-1)\big[r\log r\,I_0-(r-1)I_1\big].
\]
Let \(\sigma\) be the measure \(dt/t\) on \([1,r]\) and, for real \(s\), let
\(\ell(s)=\int t^{s+1}d\sigma\big/\!\int t^{s}d\sigma\), the mean of \(t\) under the
law proportional to \(t^s\,d\sigma\). Then \(\log\ell\) has derivative
\(\E_{s+1}\log t-\E_s\log t\), and \(s\mapsto \E_s\log t\) is strictly
increasing because its derivative is the variance of \(\log t\) under the
tilted law. Hence \(\ell\) is strictly increasing. Now
\(I_1/I_0=\ell(p-1)\) and \(\ell(-1)=\log r/(1-1/r)=r\log r/(r-1)\). As \(p>0\),
\(\ell(p-1)>\ell(-1)\), that is \((r-1)I_1>r\log r\,I_0\). The square bracket is
therefore negative, and the sign of \(\Psi_p'\) is the sign of \(-p(p-1)\).
\end{proof}

\noindent\textit{In words.} With both means held fixed, a power above one
prefers the large value to sit on as little mass as possible, and a power
below one prefers the small value to sit on as little mass as possible.

\begin{lemma}[Maximal power sum at fixed means]\label{br:lem:key}
Let \(\mu\) be a full-support law on a finite set with at least two points,
\(q=\min\mu\), \(p>0\), \(p\ne1\) and \(0<g<1\). Put
\[
S_p(\mu,g)=\sup\{\E_\mu y^p:\ y>0,\ \E_\mu y=1,\ G_\mu(y)=g\},
\qquad G_\mu(y)=\textstyle\prod_i y_i^{\mu_i}.
\]
Then the supremum is attained and
\[
S_p(\mu,g)=\begin{cases}\Psi_p(q),&p>1,\\ \Psi_p(1-q),&0<p<1.\end{cases}
\]
Every maximiser takes exactly two values. For \(p>1\), the larger value
occurs at exactly one coordinate; for \(0<p<1\), the smaller value occurs at
exactly one coordinate. In both cases that coordinate has mass \(q\).
In particular \(S_p(\mu,g)\) depends on \(\mu\) only through \(q\), and
\(S_p(\mu,g)=S_p((q,1-q),g)\). Trivially \(S_1(\mu,g)=1\), and at \(g=1\) the
only feasible vector is \(y\equiv1\).
\end{lemma}
\begin{proof}
\emph{Compactness.} On the feasible set \(F\), \(y_i\le1/\mu_i\), and
\(y_j^{\mu_j}=g\prod_{i\ne j}y_i^{-\mu_i}\ge g\prod_{i\ne j}\mu_i^{\mu_i}\).
So \(F\) is a closed subset of a compact box in \((0,\infty)^X\); the
continuous objective attains its maximum.

\emph{Stationarity.} The constraint gradients are \((\mu_i)\) and
\((g\mu_i/y_i)\). They are linearly dependent only if \(y\) is constant,
which forces \(g=1\). Hence \(F\) is a smooth submanifold of codimension two
near every point, and a maximiser satisfies
\(p\mu_iy_i^{p-1}=\kappa_1\mu_i+\kappa_2 g\mu_i/y_i\) for some multipliers
\(\kappa_1,\kappa_2\). With \(C=\kappa_2 g\), every coordinate is a positive zero of
\[
\phi(t)=pt^p-\kappa_1 t-C .
\]
As \(y\) is not constant, \(\phi\) has at least two positive zeros. If
\(p>1\), \(\phi\) is strictly convex with \(\phi(0)=-C\); a strictly convex
function with \(\phi(0)\le0\) has at most one positive zero, so \(C<0\). If
\(0<p<1\), \(\phi\) is strictly concave and the same argument gives \(C>0\).
A strictly convex or concave function has at most two zeros, so \(y\) takes
exactly two values \(x<z\).

\emph{One distinguished coordinate.} Put \(\psi(t)=\phi(t)/t\) and
\(Q(t)=p(p-1)t^p+C\), so that \(\psi'(t)=Q(t)/t^2\). If \(p>1\), \(Q\)
increases strictly from \(C<0\) to \(+\infty\); if \(p<1\), it decreases
strictly from \(C>0\) to \(-\infty\). Either way \(Q\) has one zero \(t_0\),
\(\psi\) is strictly monotone on each side of \(t_0\), and
\(\psi(x)=\psi(z)=0\) forces \(x<t_0<z\). Consequently \(Q(z)>0\) if \(p>1\),
and \(Q(x)>0\) if \(p<1\). Call that value \(v_*\) (\(z\) or \(x\)).

Suppose two coordinates \(i\ne j\) both equal \(v_*\). Take
\(h=e_i/\mu_i-e_j/\mu_j\). Then \(\sum_k\mu_kh_k=0\) and
\(\sum_k\mu_kh_k/y_k=0\), so \(h\) is tangent to \(F\). With the Lagrangian
\(L=\E_\mu y^p-\kappa_1\E_\mu y-\kappa_2 G_\mu(y)\), and
\(h^{\mathsf T}\nabla^2G_\mu h=G_\mu\big[(\sum\mu_kh_k/y_k)^2-\sum\mu_kh_k^2/y_k^2\big]\),
\[
h^{\mathsf T}\nabla^2L\,h
=\sum_k\mu_kh_k^2\Big(p(p-1)y_k^{p-2}+\frac C{y_k^2}\Big)
=\frac{Q(v_*)}{v_*^2}\sum_k\mu_kh_k^2>0 .
\]
This violates the second-order necessary condition for a local maximum on a
smooth constraint manifold. (Here \(|X|\ge3\), since the other value also
occurs, so the manifold has positive dimension.) Hence \(v_*\) occurs at
exactly one coordinate \(i\).

\emph{Which coordinate.} If \(p>1\), the maximiser is the two-valued vector
with its high value at \(i\), and its value is \(\Psi_p(\mu_i)\). If \(p<1\),
the high value has weight \(1-\mu_i\) and the value is \(\Psi_p(1-\mu_i)\). By
Lemma \ref{br:lem:weight}, both are at most the value at \(\mu_i=q\), with
equality only when \(\mu_i=q\); the vector built on a rarest atom is
feasible. For a two-point law the two-valued form is automatic and the same
comparison applies.
\end{proof}

\noindent\textit{In words.} Among positive vectors with prescribed
arithmetic and geometric means, the largest power sum is always obtained by
making one rarest coordinate different from the rest; which way it differs
depends only on whether the power is above or below one.

\subsection{The reduction}

\begin{theorem}[Block-reset reduction]\label{br:thm:reduction}
Let \(K\) be a block-reset chain \eqref{br:eq:K} with reference
\(\pi=\sum_Bw_B\mu_B\) and \(w'_C>0\) for all \(C\). For every \(c>0\),
\[
\etaKL(\pi,K)\le c\iff\etaKL(\widehat\pi,\widehat K)\le c .
\]
Consequently
\begin{equation}\label{br:eq:main}
\etaKL(\pi,K)=\etaKL(\widehat\pi,\widehat K).
\end{equation}
The coefficient depends on the block laws only through which blocks are singletons
and through their minimum masses \(q_B\). Moreover, lifting an
input \(\widehat v\) of \(\widehat K\) by placing its rare-point mass of
block \(B\) on one rarest atom of \(\mu_B\), and its other mass of block \(B\)
proportionally to \(\mu_B\) on the remaining atoms, preserves both divergences.
Hence every extremiser of \(\widehat K\) lifts to an extremiser of \(K\).
\end{theorem}

\begin{proof}
Both \(\pi\) and \(\lambda\) have full support, and so do \(\widehat\pi\)
and \(\widehat\lambda\), whose block weights are again \(w_B\) and
\(w'_C\). Lemma \ref{lem:duality} therefore applies to both chains.

Let \(f\) be real on \(X\), \(y=e^f\), and let \(y_B\) be the restriction of
\(y\) to \(X_B\). Write \(A_B=\E_{\mu_B}y_B\) and \(G_B=G_{\mu_B}(y_B)\).
Since \(\E_{\mu_C}f=\log G_C\), for \(x\in X_B\),
\[
Kf(x)=a_Bf(x)+\sum_CM_{BC}\log G_C .
\]
With \(p_B=a_B/c\ge0\) and \(\Gamma_B=\prod_CG_C^{M_{BC}/c}>0\), the dual
criterion \eqref{eq:duality} reads
\begin{equation}\label{br:eq:dual}
\Big(\sum_Bw_B\,\Gamma_B\,\E_{\mu_B}\big[y_B^{\,p_B}\big]\Big)^{c}
\le\sum_Cw'_C\,A_C\qquad\text{for all }y>0 .
\end{equation}
The right side and every \(\Gamma_B\) depend on \(y\) only through the
numbers \((A_B,G_B)_B\). The set of \(y\) with prescribed numbers is a
product over blocks. The left side is continuous and increasing in each
\(\E_{\mu_B}[y_B^{p_B}]\), because \(w_B\Gamma_B>0\). Hence
\eqref{br:eq:dual} is equivalent to the same inequality with
\(\E_{\mu_B}[y_B^{p_B}]\) replaced by
\[
\sup\{\E_{\mu_B}u^{p_B}:\ u>0,\ \E_{\mu_B}u=A_B,\ G_{\mu_B}(u)=G_B\}
=A_B^{p_B}\,S_{p_B}(\mu_B,G_B/A_B),
\]
for every admissible family \((A_B,G_B)_B\). For a block of size one,
\(G_B=A_B\) and the supremum is \(A_B^{p_B}\). For a block of size at least
two, the admissible pairs are \(0<G_B\le A_B\), and by Lemma
\ref{br:lem:key} (with the trivial cases \(p_B\in\{0,1\}\) or \(G_B=A_B\)) the
supremum depends only on \(q_B\), \(p_B\), \(A_B\) and \(G_B\). The collapsed
chain has the same weights, the same coefficients, the same admissible
pairs, and block laws with the same minima. So the reduced inequality, and
hence \eqref{br:eq:dual}, holds for \(K\) exactly when it holds for
\(\widehat K\). This proves the equivalence for each \(c>0\), and
\eqref{br:eq:main} follows by taking the infimum over admissible \(c\).

For the lift \(v\) of \(\widehat v\), on the non-rarest atoms of block \(B\)
the input \(v\) is proportional to \(\mu_B\), and so is the reference
\(\pi\). The output \(vK\) on the corresponding atoms of block \(C\) is
\(a_Cv+\big(\sum_BM_{BC}v(X_B)\big)\mu_C\), again proportional to \(\mu_C\),
as is \(\lambda\). Grouping these atoms therefore loses no divergence, and
the grouped laws are exactly \(\widehat v\), \(\widehat\pi\),
\(\widehat v\widehat K\) and \(\widehat\lambda\). Both divergences agree,
so the ratios agree; with \eqref{br:eq:main}, extremisers lift.
\end{proof}

\noindent\textit{In words.} For any chain that either stays put or jumps to a
fresh sample from some block law, the worst-case one-step entropy
contraction is the same as for a chain with at most two states per block,
and only the rarest point of each block matters.

\begin{remark}[No restriction on the laziness]\label{br:rem:norestriction}
No restriction on the laziness is needed: when some exponent
\(p_B=a_B/c\) is below one, Lemma \ref{br:lem:key} still applies, with a
single \emph{low} coordinate at a rarest atom.
This matters: for \(P_s\) below (\(s<1/4\)) one has
\(\etaKL\ge(1-2s)^2>s\), so the coefficient exceeds every stay
probability. Numerically, \(20\) of the \(24\) random chains tested by
\texttt{verification/extension\_checks.py} have a reported coefficient above
their smallest positive stay probability, and the full and collapsed
numerical estimates agree in all \(24\) to within \(10^{-15}\). The theorem does not assert that every extremiser of
\(K\) is a lift, nor does it cover cross weights that are not of the
rank-one block form \eqref{br:eq:K}; for such weights the minimum masses
alone cannot determine the coefficient.
\end{remark}

\subsection{Consequences}

Throughout, for laws \(\mu\) on a part \(A\) and \(\nu\) on a part \(B\),
and \(s\in[0,1]\), let \(P_s\) be the \emph{\(s\)-lazy bipartite reset chain}:
it stays put with probability \(s\) and otherwise samples the opposite part,
\[
P_s(a,a)=s,\ \ P_s(a,b)=(1-s)\nu_b;\qquad
P_s(b,b)=s,\ \ P_s(b,a)=(1-s)\mu_a .
\]
Its stationary law is \(\pi=(\mu/2,\nu/2)\), and \(P_{1/2}\) is the chain
\eqref{eq:weightedP}. It is block-reset with \(a_A=a_B=s\),
\(M_{AB}=M_{BA}=1-s\), \(M_{AA}=M_{BB}=0\).

\begin{corollary}[Four-state reductions]\label{br:cor:four}
For every \(s\in[0,1]\) and every \(k\ge1\), the coefficient
\(\etaKL(\pi,P_s^k)\) equals that of the collapsed chain with the same
coefficients on at most four states: four when both parts are
non-singletons, three when exactly one is a singleton, and two when both
are singletons. Each non-singleton part law \(\mu\) is replaced by
\((q_\mu,1-q_\mu)\), while each singleton is retained. In the block form,
\(P_s^k\) stays with probability \(s^k\) and otherwise resets to its own part
or to the other part. For \(s=1/2\),
\begin{equation}\label{br:eq:powers}
P_{1/2}^k=\Pi+2^{-k}(I-Q),
\end{equation}
where \(\Pi\) resets to \(\pi\) and \(Q\) resets within the current part, and
\[
\etaKL(\pi,P_{1/2}^k)\le\Big(\frac{\log(1/q)}{2\log(2/q)}\Big)^{k},
\qquad q=\min\{q_\mu,q_\nu\}.
\]
\end{corollary}
\begin{proof}
Lemma \ref{br:lem:closure} and Theorem \ref{br:thm:reduction}. For
\eqref{br:eq:powers}, write \(J\) for the cross reset, so \(J^2=Q\),
\(QJ=JQ=J\), and \(\Pi=(J+Q)/2\). Then \(P_{1/2}=(I+J)/2\) and induction gives
\(P_{1/2}^{k+1}=\tfrac12(\Pi+2^{-k}(I-Q))(I+J)
=\Pi+2^{-k-1}(I-Q)\), using \(\Pi J=\Pi\) and \((I-Q)J=0\). The bound follows
from the chain rule \(\etaKL(\pi,K_1K_2)\le\etaKL(\pi,K_1)\,\etaKL(\pi K_1,K_2)\),
stationarity of \(\pi\), and Theorem \ref{thm:bipartite}.
\end{proof}

\noindent\textit{In words.} Any laziness, and any number of steps, on a
two-part reset chain is captured exactly by a chain with at most four states.

\begin{remark}[Consistency with Theorem \ref{thm:bipartite}]\label{br:rem:thm61}
At \(s=1/2\), Corollary \ref{br:cor:four} says that Theorem
\ref{thm:bipartite} for arbitrary part laws follows from its case of
two-point parts. It is a consistency check rather than a new proof: the
four-state value still has to be computed, and Theorem \ref{thm:bipartite}
says more, namely that only \(\min\{q_\mu,q_\nu\}\) matters at this
laziness. Numerically, full chains with parts \((0.1,0.2,0.7),(0.3,0.7)\),
\((0.2,0.3,0.5),(0.15,0.25,0.6)\) and uniform \(K_{3,3}\) reproduce
\eqref{eq:bipartite} to \(10^{-15}\).
\end{remark}

\begin{remark}[The lazy star at any laziness]\label{br:cor:star}
For the star with \(n\ge3\) vertices (one singleton part, the centre, and
leaves with uniform law) and any laziness \(s\in[0,1]\), Corollary
\ref{br:cor:four} shows that the coefficient equals that of the three-state
chain with leaf law \((1/(n-1),\,1-1/(n-1))\). For \(s=1/2\) see
Corollary \ref{cor:star}, which gives the value and the extremisers.
\end{remark}

\begin{corollary}[Complete multipartite graphs]\label{br:cor:multi}
Let \(G\) be the complete multipartite graph with \(r\ge2\) parts of sizes
\(n_1,\dots,n_r\), \(N=\sum_in_i\), and let \(P\) stay put with probability
\(s\) and otherwise move to a uniform neighbour. With its stationary law
\(\pi\propto\) degree, \(\etaKL(\pi,P)\) equals the coefficient of the
block-reset chain on at most \(2r\) states with part laws
\((1/n_i,1-1/n_i)\) (a single state when \(n_i=1\)), stay probability \(s\),
reset coefficients \(M_{ij}=(1-s)n_j/(N-n_i)\) for \(j\ne i\), and block
weights \(w_i\propto n_i(N-n_i)\).
\end{corollary}
\begin{proof}
From a vertex of part \(i\), the walk moves to part \(j\ne i\) with
probability \((1-s)n_j/(N-n_i)\) and then to a uniform vertex of that part;
the degree of every vertex of part \(i\) is \(N-n_i\). Apply Theorem
\ref{br:thm:reduction}.
\end{proof}

\noindent\textit{In words.} A walk on a complete multipartite graph has the
same one-step entropy contraction as a walk on two states per part.

\subsection{Where point inputs stop being optimal}

The following two facts are proved; the window they bound is then described
numerically. For \(s\in[0,1)\) and a vertex \(x\) of conditional mass
\(\rho_x\), the point ratio of \(P_s\) is
\[
c_0(\rho_x,s)=\frac{s\log(2s/\rho_x)+(1-s)\log(2(1-s))}{\log(2/\rho_x)}
=s-\frac{h(s)-(1-s)\log2}{\log(2/\rho_x)},
\]
with \(h(s)=-s\log s-(1-s)\log(1-s)\). For \(s\in[1/2,1)\),
\(h(s)\ge2(1-s)\log2>(1-s)\log2\) by concavity of \(h\), so the largest
point ratio is at a vertex of the smallest conditional mass
\(q=\min\{q_\mu,q_\nu\}\).

\begin{proposition}[Two obstructions to point optimality]\label{br:prop:edges}
\begin{enumerate}[label=(\roman*)]
\item Assume the part containing a rarest vertex has at least two vertices.
If \(1/2\le s<1\) and \(c_0(q,s)>(1-s)(1-q)\), then
\(\etaKL(\pi,P_s)>c_0(q,s)\); no point input is an extremiser.
\item If \((1-2s)^2>\max_xc_0(\rho_x,s)\), then
\(\etaKL(\pi,P_s)\ge(1-2s)^2\) exceeds every point ratio, so again no point
input is an extremiser.
\end{enumerate}
\end{proposition}
\begin{proof}
(i) Let \(a\) be a rarest vertex, in part \(A\), and \(b\) any vertex of the
other part. Put \(v_\epsilon=(1-\epsilon)\delta_a+\epsilon\delta_b\). Its
output gives mass \(\epsilon(1-s)\mu_j\) to each \(j\in A\setminus\{a\}\),
which the point output does not charge; every other output coordinate is a
smooth function of \(\epsilon\) and positive at \(\epsilon=0\). Hence
\[
\KL(v_\epsilon P_s\|\pi)=\KL(\delta_aP_s\|\pi)
+(1-s)(1-q)\,\epsilon\log\epsilon+O(\epsilon),\qquad
\KL(v_\epsilon\|\pi)=\log(2/q)+\epsilon\log\epsilon+O(\epsilon).
\]
The output divergence minus \(c_0(q,s)\) times the input divergence is
\(\big((1-s)(1-q)-c_0(q,s)\big)\epsilon\log\epsilon+O(\epsilon)\), which is
positive for small \(\epsilon>0\) under the hypothesis. Since \(c_0(q,s)\)
is the largest point ratio, no point input is optimal. This is the argument
of \eqref{eq:perturb}.

(ii) The chain is reversible with respect to \(\pi\), and
\(\phi=\one_A-\one_B\) has \(\E_\pi\phi=0\) and \(P_s\phi=(2s-1)\phi\).
For \(v_t=\pi(1+t\phi)\), reversibility gives
\(v_tP_s=\pi(1+t(2s-1)\phi)\). Expanding both divergences to second order,
\(\KL(v_tP_s\|\pi)/\KL(v_t\|\pi)\to(2s-1)^2\) as \(t\to0\).
\end{proof}

\noindent\textit{In words.} Laziness a little above one half, or well below
it, makes some spread-out starting law worse than every single vertex.

For uniform parts, \(q=1/n\), the root \(s_+(n)\) of
\(c_0(q,s)=(1-s)(1-q)\) above \(1/2\) is \(1/2\) for \(n=2\) and
\(0.52054\), \(0.52885\), \(0.53307\), \(0.53877\), \(0.53228\) for
\(n=3,4,5,10,100\) (computed by root finding); the crossing \(s_-(n)\) of \(c_0\) with \((1-2s)^2\) is
\(0.26836\), \(0.26204\), \(0.25967\), \(0.25839\), \(0.25595\), \(0.25302\)
for \(n=2,3,4,5,10,100\). The upper edge \(s_+(n)\) is not monotone in
\(n\): it peaks near \(n=10\) and then decreases towards \(1/2\).
At \(n=3\) and \(220\)-digit precision, the perturbation in (i) with
\(\epsilon=10^{-200}\) gains \(1.3\times10^{-200}\) at \(s=0.525\) and loses
\(2.1\times10^{-200}\) at \(s=0.515\), as (i) predicts.

\begin{remark}[Numerical observations, not proved]\label{br:rem:numerics}
We report the largest value among point-input evaluations, multistart
numerical searches, and the local spectral lower bound. Accepted
optimisation candidates are reevaluated at \(40\)-digit working precision.
Local spectral values may be limiting values rather than values attained at
a nonreference input. These calculations provide numerical evidence, without
certified error bounds or certified global optimality. They were produced by
\texttt{verification/extension\_checks.py} and cross-checked on the full
chains.
\begin{enumerate}[label=(\alph*)]
\item \emph{Laziness window on \(K_{n,n}\)} (\(n\in\{2,3,4,5,10,100\}\),
\(s\in[0.20,0.70]\) in steps of \(0.01\)). For \(s<1/2\) the reported numerical estimate
agrees with \(\max\{c_0(1/n,s),(1-2s)^2\}\) to \(10^{-15}\) at every grid
point; no improvement over the point ratio larger than \(10^{-9}\) is found at
any grid point from the first one above \(s_-(n)\) up to \(1/2\).
Above \(1/2\), the first grid point with an improvement above \(10^{-9}\)
is \(0.51,0.53,0.54,0.55,0.55,0.55\) for \(n=2,3,4,5,10,100\), and every
later grid point also shows one. For \(n=2,3\) this is the first grid point
above \(s_+(n)\), consistent with point optimality on exactly
\([s_-(n),s_+(n)]\). For larger \(n\) it lies one or two grid steps above
\(s_+(n)\): just above \(s_+(n)\) the gain from (i) is of
\(\epsilon\log\epsilon\) type with extremely small \(\epsilon\) and is
invisible in double precision, so there the grid detects failure only from
outside and Proposition \ref{br:prop:edges}(i) is the sharper statement. For
\(n=2\) the grid values above \(1/2\) agree with \(s^2\).
Full six- and eight-state chains agree with the four-state values to
\(10^{-15}\), for example \(0.340916\) versus point \(0.340026\) at
\((n,s)=(3,0.55)\) and \(0.389576\) versus \(0.379126\) at \((3,0.6)\).
\item \emph{Powers at \(s=1/2\).} Point inputs are not optimal for
\(P_{1/2}^2\) and \(P_{1/2}^3\):
\(\etaKL\approx0.0721278\) (point \(0.0644755\)) and \(0.0185000\)
(point \(0.0164340\)) for \(n=3\); \(0.0886389\) (point \(0.0829304\)) and
\(0.0242136\) (point \(0.0227244\)) for \(n=5\). The best input found on four states
for \(n=3,k=2\) lifts to a full-support input of \(K_{3,3}\). All values are
below the bound in Corollary \ref{br:cor:four}.
\item \emph{Complete multipartite graphs at \(s=1/2\).} Unlike the bipartite
case, point inputs fail to be optimal in every tested example:
\(K_{2,3,4}\): \(0.310885\) versus point \(0.308287\);
\(K_{3,3,3}\): \(0.286790\) versus \(0.276803\);
\(K_{1,2,3}\): \(0.288081\) versus \(0.285553\);
\(K_{2,2,2,3}\): \(0.289083\) versus \(0.276421\).
Full and collapsed chains agree to \(10^{-15}\).
\end{enumerate}
An analytic description of the window \([s_-(n),s_+(n)]\), and closed forms
for powers and for multipartite graphs, are left open. The reductions above
turn each into a fixed small-dimensional problem.
\end{remark}
